\begin{proof}[Proof of \cref{obj:thm:original-record-attainable-rate}]
\begin{enumerate}
\item Fix an admissible tuple \(v=(p,\alpha,\beta,\gamma)\). Its parameter restrictions give
\[
1<p\le2,\qquad
0<\alpha,\beta\le1,\qquad
\frac14\le\gamma\le1.
\]
Using the quantities in the statement, write
\[
D_p=1+2\alpha+\frac{\beta}{q}+\frac{S}{2\gamma}.
\]
Then
\[
0<q\le\frac12,\qquad S>0,\qquad t\ge0,\qquad D_p>1,
\qquad E_0>0,\qquad E_4>0.
\]
In particular,
\[
0<E\le\gamma\le1,\qquad
\frac{E}{S}<2,\qquad
\frac{E}{p-1}<1.
\]
Indeed, the three upper bounds follow from
\[
\frac{E}{\gamma}
\le\frac{2q}{2\gamma+q}\le1,
\qquad
\frac{E}{S}\le\frac{2}{D_p}<2,
\qquad
\frac{E}{p-1}
\le\frac{2\gamma}{p(2\gamma+q)}<1.
\]
The first uses \(q\le1/2\le2\gamma\), and the last uses \(p>1\) and \(q>0\).

For every positive argument of the geometric-series factor,
\[
0<2^{-x}<1,\qquad A(x)>0.
\]
Since \(\alpha>0\) and \(\lambda_0>0\), it follows that \(B_*\ge0\). Thus
\[
C^{\mathrm{att}}_v
=16\{16+5B_*+1024\sqrt{A_*}\}
\ge256>0.
\]
% lean: CausalSmith.Stat.FinitepHomogeneityDensegamma.CAtt_ge_256

\item The rule named \(\psi_{n,v}\) in
\cref{obj:thm:whole-null-calibration-resolved} is precisely
\(\phi^*_{n,v}\). That result supplies, for every integer \(n\ge2\),
\[
\mathsf R_n(P,\phi^*_{n,v})\le0.0075
\qquad(P\in H_0(v)).
\]
It also supplies power, for \(n\ge4\), at the calibrated separation
\[
R_{\mathrm{att}}(n,v)
=\frac{16}{3}
\left\{
16M^{-\gamma}+5B+1024M^{1/4}\sqrt{\Lambda_{n,v}}
\right\},
\qquad B=\mathfrak b_{n,v},
\]
provided \(R_{\mathrm{att}}(n,v)<D_v\).

The smallest sample sizes can be handled directly.
By \cref{obj:lem:causal-nonempty},
\[
d_0\le D_v\le\frac12.
\]
For \(n\in\{2,3\}\), the bound \(E\le1\) gives
\[
n^{-E}\ge n^{-1}\ge\frac13,
\qquad
D_v\le\frac12<\frac{256}{3}
\le C^{\mathrm{att}}_v n^{-E}.
\]
Consequently,
\[
C^{\mathrm{att}}_v n^{-E}<D_v
\quad\Longrightarrow\quad n\ge4
\qquad(n\ge2).
\]
The zero rule at \(n=2,3\) has the asserted size by the same calibration result.
We next establish the deterministic estimates needed to compare
\(R_{\mathrm{att}}(n,v)\) with the stated public level.
% lean: CausalSmith.Stat.FinitepHomogeneityDensegamma.attainment_small_sample_saturated

\item Let \(n\ge4\), and use the tuning of
\cref{obj:def:explicit-score-ledger,obj:def:blocks}. Explicitly,
\[
s=\lfloor n/2\rfloor\ge2,\qquad
a=s^{-E}\in(0,1],\qquad
M=\mathcal R(a^{-1/\gamma}),\qquad
K=\mathcal R\!\left(\max\{M,a^{-1/S}\}\right).
\]
Upper dyadic rounding obeys
\[
x_{\mathrm{dyad}}\le\mathcal R(x_{\mathrm{dyad}})
\le2x_{\mathrm{dyad}}
\qquad(x_{\mathrm{dyad}}\ge1),
\qquad
\mathcal R(2^j)=2^j\qquad(j\in\mathbb N_0),
\]
by the lower and upper ceiling bounds for \(\log_2 x_{\mathrm{dyad}}\).
It is also increasing. Hence
\[
a^{-1/\gamma}\le M\le2a^{-1/\gamma},
\qquad
K\ge a^{-1/S},
\]
and therefore
\[
M^{-\gamma}\le a,\qquad K^{-S}\le a.
\]

The bounds \(E/\gamma\le1\), \(E/S<2\), and \(E/(p-1)<1\) give
\[
a^{-1/\gamma}=s^{E/\gamma}\le s,
\qquad
a^{-1/S}=s^{E/S}\le s^2,
\qquad
1\le a^{-1/(p-1)}=s^{E/(p-1)}\le s.
\]
Thus \(M\le2s\). For the fine rank, if \(a^{-1/S}\le M\), dyadic idempotence gives \(K=M\le2s\le2s^2\); if \(M<a^{-1/S}\), rounding gives \(K\le2a^{-1/S}\le2s^2\). In either case,
\[
M\le2s,\qquad K\le2s^2.
\]

For the increment-cutoff bounds, the public quantities in either score branch are
\[
L=\log_2(K/M)\in\mathbb N_0,\qquad
R_j=2^jM\quad(0\le j\le L),
\]
\[
T_j=
\mathcal R\!\left(
\max\left\{
1,\min\left\{
a^{-1/(p-1)},
\left[
\frac{R_j^{-\alpha}}{a(R_j/K)^{\lambda_0}}
\right]^{1/(p-1)}
\right\}
\right\}
\right)
\quad(1\le j\le L).
\]
The single-cutoff score uses the main cutoff
\[
T_0=\mathcal R(a^{-1/(p-1)}).
\]
Since the argument defining each \(T_j\) lies between \(1\) and
\(a^{-1/(p-1)}\), rounding and monotonicity yield
\[
1\le T_0\le2s,\qquad
1\le T_j\le T_0\quad(1\le j\le L).
\]
These statements also cover \(L=0\), when the increment index set is empty.
% lean: CausalSmith.Stat.FinitepHomogeneityDensegamma.ledger_rate_bounds

\item We now bound the bias ledger. The main-cutoff estimate is
\[
T_0^{1-p}
\le\left(a^{-1/(p-1)}\right)^{1-p}
=a.
\]
In the single-cutoff branch, this and \(K^{-S}\le a\) immediately give
\[
B=400K^{-S}+20T_0^{1-p}\le420a\le B_*a.
\]

For the multilevel branch, \(S<\gamma\) and \(\alpha<q\). Define the positive local tail targets by
\[
z_{\mathrm{tail},j}
=\frac{R_j^{-\alpha}}{a(R_j/K)^{\lambda_0}}>0
\qquad(1\le j\le L).
\]
The cutoff formula gives
\[
T_j\ge
\min\left\{a^{-1/(p-1)},z_{\mathrm{tail},j}^{1/(p-1)}\right\}.
\]
If the minimum is the first term, its \((1-p)\)-power is \(a\); if it is the second term, that power is
\[
\left(z_{\mathrm{tail},j}^{1/(p-1)}\right)^{1-p}
=z_{\mathrm{tail},j}^{-1}
=a(R_j/K)^{\lambda_0}R_j^\alpha.
\]
Since \(1-p<0\), both cases imply
\[
T_j^{1-p}
\le a+a(R_j/K)^{\lambda_0}R_j^\alpha.
\]

The increment ranks satisfy \(M\ge1\) and \(R_L\le K\). Forward and backward geometric summation therefore give, for every \(r_{\mathrm{geom}}>0\),
\[
\begin{aligned}
\sum_{j=1}^{L}R_{j-1}^{-r_{\mathrm{geom}}}
&\le\sum_{j=0}^{L-1}2^{-jr_{\mathrm{geom}}}
\le A(r_{\mathrm{geom}}),\\
\sum_{j=1}^{L}(R_j/K)^{r_{\mathrm{geom}}}
&\le\sum_{j=0}^{L-1}2^{-jr_{\mathrm{geom}}}
\le A(r_{\mathrm{geom}}),\\
\sum_{j=1}^{L}R_j^{r_{\mathrm{geom}}}
&\le K^{r_{\mathrm{geom}}}A(r_{\mathrm{geom}}).
\end{aligned}
\]
For the second inequality, use \(R_j/K\le2^{-(L-j)}\); the third follows by factoring out \(K^{r_{\mathrm{geom}}}\). Each bound remains valid for an empty sum.

Using \(a_j=40R_{j-1}^{-\alpha}\), \(R_j=2R_{j-1}\), and \(2^\alpha\le2\), we obtain
\[
\begin{aligned}
a_jT_j^{1-p}
&\le40aR_{j-1}^{-\alpha}
 +40a\,2^\alpha(R_j/K)^{\lambda_0}\\
&\le80a\left\{R_{j-1}^{-\alpha}+(R_j/K)^{\lambda_0}\right\}.
\end{aligned}
\]
Summation consequently gives
\[
\sum_{j=1}^{L}a_jT_j^{1-p}
\le80a\{A(\alpha)+A(\lambda_0)\}.
\]
Substituting into the multilevel bias ledger proves
\[
\begin{aligned}
B
&=400K^{-S}+20T_0^{1-p}
  +10\sum_{j=1}^{L}a_jT_j^{1-p}\\
&\le\left[420+800\{A(\alpha)+A(\lambda_0)\}\right]a
=B_*a.
\end{aligned}
\]
Thus this bias estimate holds in both score branches.
% lean: CausalSmith.Stat.FinitepHomogeneityDensegamma.ledger_bias_bound

\item Consider first the single-cutoff covariance branch,
\[
S\ge\gamma\quad\text{or}\quad\alpha\ge q.
\]
If \(S\ge\gamma\), then \(E\le\gamma\le S\). If \(\alpha\ge q\), then
\[
E\le E_0\le q\le\alpha<S.
\]
Thus \(a^{-1/S}=s^{E/S}\le s\). The same two cases for fine-rank rounding used above now give
\[
K\le2s.
\]

Since \(0\le2-p\le1\) and \(T_0\le2a^{-1/(p-1)}\), the main second-moment quantity satisfies
\[
V_0=32T_0^{2-p}
\le32\,2^{2-p}a^{-t}
\le64a^{-t}.
\]
The single-cutoff formulas \(L_1=16\sqrt{V_0}\) and \(L_2=4\sqrt{KV_0}\) yield the exact identity
\[
\Lambda_{n,v}
=\frac{256V_0}{s}+\frac{32KV_0}{s(s-1)}.
\]
Because \(s\ge2\) and \(K\le2s\),
\[
32K\le64s\le128(s-1).
\]
It follows that
\[
\Lambda_{n,v}
\le\frac{384V_0}{s}
\le\frac{24576a^{-t}}{s}.
\]

The tail-channel identity and the choice \(E\le E_0\) give
\[
E_0\left(2+t+\frac{1}{2\gamma}\right)=1,
\qquad
E\left(2+t+\frac{1}{2\gamma}\right)\le1.
\]
Hence
\[
\frac{a^{-1/(2\gamma)}a^{-t}}{s}
=s^{E(1/(2\gamma)+t)-1}
\le s^{-2E}=a^2.
\]
Finally, \(M\le2a^{-1/\gamma}\) implies
\[
\begin{aligned}
\sqrt M\,\Lambda_{n,v}
&\le24576\sqrt2\,
 \frac{a^{-1/(2\gamma)}a^{-t}}{s}\\
&\le24576\sqrt2\,a^2
\le A_*a^2,
\end{aligned}
\]
where the last inequality follows from the displayed definition of \(A_*\).
% lean: CausalSmith.Stat.FinitepHomogeneityDensegamma.ledger_single_covariance_rate

\item Continue with the multilevel branch \(S<\gamma\), \(\alpha<q\).
Here \(a^{-1/\gamma}\le a^{-1/S}\). If \(a^{-1/S}\le M\), then
\[
K=M\le2a^{-1/\gamma}\le2a^{-1/S};
\]
if \(M<a^{-1/S}\), rounding gives the same final bound. Thus
\[
K\le2a^{-1/S},
\qquad
aK^\alpha
\le2^\alpha a^{1-\alpha/S}
=2^\alpha a^{\beta/S}\le2.
\]

The refinement depth satisfies
\[
2^L\le R_L\le K\le2a^{-1/S},
\qquad
L\log2\le\log2-\frac{\log a}{S}.
\]
The exponential tangent inequality, applied at \(-\log a-1\), gives
\[
-\log a
\le\exp(-\log a-1)=a^{-1}\exp(-1),
\qquad
-a\log a\le\exp(-1).
\]
Consequently,
\[
aL
\le a+\frac{-a\log a}{S\log2}
\le1+\frac{1}{\exp(1)S\log2}
=D_*.
\]

The bound \(T_j^{1-p}\le a+a(R_j/K)^{\lambda_0}R_j^\alpha\) can be rewritten using
\[
a(R_j/K)^{\lambda_0}R_j^\alpha
=aK^\alpha(R_j/K)^{\alpha+\lambda_0}.
\]
Geometric summation, \(aK^\alpha\le2\), and the fact that \(A\) decreases on positive arguments now give
\[
\begin{aligned}
\sum_{j=1}^{L}T_j^{1-p}
&\le aL+aK^\alpha A(\alpha+\lambda_0)\\
&\le D_*+2A(\alpha).
\end{aligned}
\]
The same geometric estimates supply
\[
\sum_{j=1}^{L}a_j\le80A(\alpha),
\]
and, from \(d_j=110R_{j-1}^{-s_0}+20T_j^{1-p}\),
\[
\sum_{j=1}^{L}d_j
\le110A(s_0)+20D_*+40A(\alpha).
\]
Moreover, \(T_j\le T_0\) and \(2-p\ge0\) imply
\[
\sqrt{V_j}\le\sqrt{V_0},
\qquad
\sum_{j=1}^{L}a_j\sqrt{V_j}
\le80A(\alpha)\sqrt{V_0}.
\]
Since \(T_0\ge1\), we also have \(\sqrt{V_0}\ge1\).
Inserting these bounds into the singleton ledger gives
\[
\begin{aligned}
L_1
&=16\sqrt{V_0}
 +8\sum_{j=1}^{L}(d_j+a_j\sqrt{V_j})\\
&\le
\left\{16+880A(s_0)+160D_*+960A(\alpha)\right\}
\sqrt{V_0}
=L_*\sqrt{V_0}.
\end{aligned}
\]
Using \(V_0\le64a^{-t}\), we obtain
\[
L_1^2\le64L_*^2a^{-t}.
\]
Together with \(a^{-1/(2\gamma)}a^{-t}/s\le a^2\), this proves
\[
\sqrt M\,\frac{L_1^2}{s}
\le\sqrt2\,64L_*^2
\frac{a^{-1/(2\gamma)}a^{-t}}{s}
\le\sqrt2\,64L_*^2a^2.
\]
% lean: CausalSmith.Stat.FinitepHomogeneityDensegamma.ledger_multilevel_singleton_covariance_rate

\item Still in the multilevel branch, define the backward summability exponent by
\[
\eta=1-(\alpha+\lambda_0)t.
\]
Because \(\alpha<q\) and \(t\ge0\),
\[
\eta\ge1-(q+\lambda_0)t
=\frac{(p-1)(p+6)}{4p}
=\eta_0>0.
\]
Also,
\[
0\le\alpha t=1-\eta-\lambda_0t<1.
\]

The positive-power cutoff estimate follows from
\[
T_j\le2\max\{1,z_{\mathrm{tail},j}^{1/(p-1)}\}.
\]
Indeed, using \(2^{2-p}\le2\) gives
\[
V_j=32T_j^{2-p}
\le64\max\{1,z_{\mathrm{tail},j}^{t}\},
\]
and therefore
\[
\sqrt{V_j}
\le\sqrt{64}\left(1+z_{\mathrm{tail},j}^{t/2}\right).
\]
Multiplying by \(\sqrt{R_j}\) and expanding the local target yields
\[
\sqrt{R_jV_j}
\le\sqrt{64}
\left\{
\sqrt{R_j}
+a^{-t/2}K^{\lambda_0t/2}R_j^{\eta/2}
\right\}.
\]
Backward geometric summation implies
\[
\sum_{j=1}^{L}\sqrt{R_j}
\le\sqrt K\,A(1/2),
\qquad
\sum_{j=1}^{L}R_j^{\eta/2}
\le K^{\eta/2}A(\eta/2)
\le K^{\eta/2}A(\eta_0/2).
\]
The fine-rank amplitude estimate also gives
\[
\begin{aligned}
\sqrt K
&=(aK^\alpha)^{t/2}
  a^{-t/2}K^{(1-\alpha t)/2}\\
&\le2^{t/2}a^{-t/2}K^{(1-\alpha t)/2}.
\end{aligned}
\]
Since \(\lambda_0t+\eta=1-\alpha t\), summing the guarded increment bound proves
\[
\begin{aligned}
\sum_{j=1}^{L}\sqrt{R_jV_j}
&\le\sqrt{64}
 \left\{2^{t/2}A(1/2)+A(\eta_0/2)\right\}
 a^{-t/2}K^{(1-\alpha t)/2}\\
&=G_*a^{-t/2}K^{(1-\alpha t)/2}.
\end{aligned}
\]
Now
\[
0<\frac{1-\alpha t}{2}\le\frac12,
\qquad
t+\frac{1-\alpha t}{S}=\frac{1+\beta t}{S}.
\]
Thus \(K\le2a^{-1/S}\) converts the last bound into
\[
\sum_{j=1}^{L}\sqrt{R_jV_j}
\le\sqrt2\,G_*a^{-(1+\beta t)/(2S)}.
\]

Using the definition of \(L_2\) and \((x+y)^2\le2x^2+2y^2\) for its two nonnegative summands, we get
\[
\begin{aligned}
L_2^2
&=16\left(
\sqrt{MV_0}+\sum_{j=1}^{L}\sqrt{R_jV_j}
\right)^2\\
&\le32MV_0
 +32\left(\sum_{j=1}^{L}\sqrt{R_jV_j}\right)^2\\
&\le4096s\,a^{-t}
 +64G_*^2a^{-(1+\beta t)/S}.
\end{aligned}
\]
Here the last line uses \(M\le2s\) and \(V_0\le64a^{-t}\).
Since \(s(s-1)\ge s^2/2\), it follows that
\[
\frac{2L_2^2}{s(s-1)}
\le\frac{16384a^{-t}}{s}
 +\frac{256G_*^2a^{-(1+\beta t)/S}}{s^2}.
\]

For the second noise channel, direct algebra gives
\[
2+t=\frac1q,
\qquad
2S+1+\beta t+\frac{S}{2\gamma}=D_p,
\]
and hence
\[
E_4\left(2+\frac{1+\beta t}{S}+\frac{1}{2\gamma}\right)=2.
\]
Using \(E\le E_4\), we obtain
\[
\frac{a^{-1/(2\gamma)}a^{-(1+\beta t)/S}}{s^2}
=s^{E(1/(2\gamma)+(1+\beta t)/S)-2}
\le s^{-2E}=a^2.
\]
Combining this with the tail-channel balance gives
\[
\sqrt M\,\frac{2L_2^2}{s(s-1)}
\le\sqrt2\{16384+256G_*^2\}a^2.
\]
Adding the bound \(\sqrt M\,L_1^2/s\le\sqrt2\,64L_*^2a^2\) therefore proves
\[
\begin{aligned}
\sqrt M\,\Lambda_{n,v}
&=\sqrt M\left\{\frac{L_1^2}{s}
                 +\frac{2L_2^2}{s(s-1)}\right\}\\
&\le\sqrt2\{64L_*^2+16384+256G_*^2\}a^2\\
&\le A_*a^2.
\end{aligned}
\]
This includes \(L=0\), since all the increment sums then vanish. Together with the single-cutoff bound \(\sqrt M\,\Lambda_{n,v}\le24576\sqrt2\,a^2\le A_*a^2\), the covariance estimate holds in both branches.
% lean: CausalSmith.Stat.FinitepHomogeneityDensegamma.ledger_multilevel_covariance_rate_of_singleton

\item The preceding estimates now give
\[
M^{-\gamma}\le a,\qquad
B\le B_*a,\qquad
M^{1/4}\sqrt{\Lambda_{n,v}}
=\sqrt{\sqrt M\,\Lambda_{n,v}}
\le\sqrt{A_*}\,a.
\]
Consequently,
\[
R_{\mathrm{att}}(n,v)
\le\frac{16}{3}
\{16+5B_*+1024\sqrt{A_*}\}\,a.
\]
The relation \(n\le3s\), together with \(0<E\le1\), implies
\[
a=s^{-E}\le3^E n^{-E}\le3n^{-E}.
\]
Thus
\[
R_{\mathrm{att}}(n,v)\le C^{\mathrm{att}}_v n^{-E}.
\]

If \(C^{\mathrm{att}}_v n^{-E}<D_v\), the bound \(D_v\le C^{\mathrm{att}}_v n^{-E}\) for \(n\in\{2,3\}\) ensures \(n\ge4\). The last inequality then gives
\[
R_{\mathrm{att}}(n,v)<D_v,\qquad
d(P)\ge C^{\mathrm{att}}_v n^{-E}
\ \Longrightarrow\
d(P)\ge R_{\mathrm{att}}(n,v).
\]
The power assertion of
\cref{obj:thm:whole-null-calibration-resolved} therefore yields
\[
1-\mathsf R_n(P,\phi^*_{n,v})\le0.0075
\]
for every stipulated alternative.

To obtain the capped-radius bound, observe that
\(C^{\mathrm{att}}_v n^{-E}>0\) and that both error bounds are below \(1/10\).
When \(C^{\mathrm{att}}_v n^{-E}<D_v\), the explicit test is admissible in
\cref{obj:def:testing-risk}, so
\[
B_n\!\left(v,C^{\mathrm{att}}_v n^{-E}\right)
\le0.0075\le\frac1{10}.
\]
For an empty alternative set its worst-error supremum has value zero, giving the same inequality. The defining set in
\cref{obj:def:critical-radius} therefore contains
\(C^{\mathrm{att}}_v n^{-E}\).
When \(D_v\le C^{\mathrm{att}}_v n^{-E}\), its sentinel \(D_v\) instead gives
\[
r_n^*(v)\le D_v\le C^{\mathrm{att}}_v n^{-E}.
\]
The defining set is nonempty and bounded below by zero in both cases. Hence
\[
r_n^*(v)\le C^{\mathrm{att}}_v n^{-E}
\qquad(n\ge2).
\]
% lean: CausalSmith.Stat.FinitepHomogeneityDensegamma.original_record_attainment_finite

\item Let \(n\ge N^{\mathrm{att}}_v\). By the ceiling definition,
\[
n\ge
\max\left\{4,
\left(\frac{2C^{\mathrm{att}}_v}{d_0}\right)^{1/E}\right\}.
\]
All bases are positive, and \(-E<0\). Taking this negative power therefore gives
\[
n^{-E}
\le
\left[
\left(\frac{2C^{\mathrm{att}}_v}{d_0}\right)^{1/E}
\right]^{-E}
=\left(\frac{2C^{\mathrm{att}}_v}{d_0}\right)^{-1}.
\]
Thus
\[
C^{\mathrm{att}}_v n^{-E}
\le\frac{d_0}{2}<d_0,
\]
which proves the witness-threshold assertion. The bounds \(M\le2s\), \(K\le2s^2\), \(1\le T_0\le2s\), and \(1\le T_j\le T_0\) for \(1\le j\le L\) give the tuning assertions.
% lean: CausalSmith.Stat.FinitepHomogeneityDensegamma.attainment_threshold

\item For compact uniformity, let
\[
\mathcal K\subseteq\mathcal V
\]
be compact in the product topology. At every admissible tuple,
\[
p-1>0,\quad S>0,\quad\gamma>0,\quad q>0,\quad D_p>0,
\quad s_0>0,\quad\lambda_0>0,\quad\eta_0>0.
\]
In particular, every geometric denominator appearing in the constants satisfies
\[
1-2^{-x}>0
\qquad
\left(
x\in
\{\alpha,\lambda_0,s_0,1/2,\eta_0/2\}
\right).
\]
The displayed formulas consequently show that \(E(v)\) and
\(C^{\mathrm{att}}_v\) are continuous at every admissible tuple: their operations are sums, products, quotients with nonzero denominators, minima, square roots, and real powers with positive bases.

Define the real threshold target on the admissible domain by
\[
\Xi(v)
=\max\left\{4,
\left(\frac{2C^{\mathrm{att}}_v}{d_0}\right)^{1/E(v)}
\right\}
\qquad(v\in\mathcal V).
\]
Since \(E(v)>0\), \(C^{\mathrm{att}}_v>0\), and \(d_0>0\), this function is continuous there as well.

If \(\mathcal K\) is nonempty, compactness supplies tuples
\(v_C,v_E,v_N\in\mathcal K\) satisfying
\[
C^{\mathrm{att}}_{v_C}
=\max_{v\in\mathcal K}C^{\mathrm{att}}_v,
\qquad
E(v_E)=\min_{v\in\mathcal K}E(v),
\qquad
\Xi(v_N)=\max_{v\in\mathcal K}\Xi(v).
\]
Choose
\[
C_{\mathcal K}=C^{\mathrm{att}}_{v_C},
\qquad
N_{\mathcal K}=\Xi(v_N)+1,
\qquad
e_{\mathcal K}=E(v_E)>0.
\]
The ceiling bound gives, for every \(v\in\mathcal K\),
\[
C^{\mathrm{att}}_v\le C_{\mathcal K},
\qquad
N^{\mathrm{att}}_v=\lceil\Xi(v)\rceil
<\Xi(v)+1\le N_{\mathcal K},
\qquad
e_{\mathcal K}\le E(v).
\]
If \(\mathcal K\) is empty, choose
\[
(C_{\mathcal K},N_{\mathcal K},e_{\mathcal K})=(0,0,1);
\]
the universal bounds are then vacuous.
% lean: CausalSmith.Stat.FinitepHomogeneityDensegamma.attainment_compact_uniformity

\item Finally, let \(w=(\alpha,\beta,\gamma)\in\mathcal W\), and define its matched admissible tuple by
\[
v_{\mathrm b}=(2,\alpha,\beta,\gamma)\in\mathcal V.
\]
At this tuple,
\[
q=\frac12,\qquad t=0,\qquad S=\alpha+\beta,
\]
so substitution into the two exponent formulas gives
\[
E_0(v_{\mathrm b})=\frac{2\gamma}{4\gamma+1},
\qquad
E_4(v_{\mathrm b})
=\frac{2S}{1+2S+S/(2\gamma)}.
\]
Taking their minimum proves
\[
E(v_{\mathrm b})=E_{\mathrm b}.
\]

The definitions in
\cref{obj:def:bounded-model,obj:def:bounded-null,obj:def:sharp-frontier-scales}
give
\[
\mathcal M_w^{\mathrm b}\subseteq\mathcal M_{v_{\mathrm b}},
\qquad
H_0^{\mathrm b}(w)\subseteq H_0(v_{\mathrm b}),
\qquad
\phi^{*,\mathrm b}_{n,w}=\phi^*_{n,v_{\mathrm b}},
\]
and
\[
C^{\mathrm{att}}_{(2,w)}n^{-E_{\mathrm b}}
=C^{\mathrm{att}}_{v_{\mathrm b}}n^{-E(v_{\mathrm b})}>0.
\]
Moreover, \cref{obj:prop:bounded-submodel-comparison} supplies the exact cap identity
\[
D_w^{\mathrm b}=D_{v_{\mathrm b}}\ge d_0.
\]
The original-model size bound therefore applies to every bounded null law. If
\[
C^{\mathrm{att}}_{(2,w)}n^{-E_{\mathrm b}}<D_w^{\mathrm b},
\]
the cap identity makes the original-model power condition hold at
\(v_{\mathrm b}\). Thus every bounded-model law with distance at least this level satisfies
\[
1-\mathsf R_n(P,\phi^{*,\mathrm b}_{n,w})\le0.0075.
\]

When this positive level is below \(D_w^{\mathrm b}\), the resulting bounded-model risk bound is
\[
B_n^{\mathrm b}
\!\left(w,C^{\mathrm{att}}_{(2,w)}n^{-E_{\mathrm b}}\right)
\le0.0075\le\frac1{10},
\]
so the level belongs to the defining set of
\cref{obj:def:bounded-radius}. When it is at least \(D_w^{\mathrm b}\), the sentinel gives the radius bound directly. Therefore, for every \(n\ge2\),
\[
r_n^{*,\mathrm b}(w)
\le C^{\mathrm{att}}_{(2,w)}n^{-E_{\mathrm b}}.
\]
% lean: CausalSmith.Stat.FinitepHomogeneityDensegamma.bounded_record_attainment_finite
\end{enumerate}
\end{proof}
